\clearpage
\renewcommand{\bibname}{参考文献}
\begin{thebibliography}{99}
\addcontentsline{toc}{chapter}{参考文献}
\zihao{-4}
\setlength{\itemsep}{2pt}
\setlength{\parskip}{0pt}

\bibitem{lewis2020}
LEWIS P, PEREZ E, PIKTUS A, et al. Retrieval-Augmented Generation for Knowledge-Intensive NLP Tasks[C]//Advances in Neural Information Processing Systems 33. Red Hook, NY: Curran Associates, 2020: 9459-9474.

\bibitem{baijing2026}
白晶, 张艳明, 扈韵绮, 等. 基于大语言模型和RAG技术的高校知识库智能问答系统构建与评价[J]. 科技视界, 2026(3): 110-114.

\bibitem{yuan2026}
袁乐, 刘绍华, 王禹, 等. 大语言模型检索增强生成优化技术研究综述[J]. 计算机学报, 2026, 49(2): 383-422. DOI:10.11897/SP.J.1016.2026.00383.

\bibitem{zhang2026}
张犬俊, 谢杨, 房春荣, 等. 检索增强生成在软件工程中的应用综述[J]. 软件学报, 2026, 37(3): 1316-1339. DOI:10.13328/j.cnki.jos.007567.

\bibitem{asai2024}
ASAI A, WU Z, WANG Y, et al. Self-RAG: Learning to Retrieve, Generate, and Critique through Self-Reflection[C]//The Twelfth International Conference on Learning Representations. 2024.

\bibitem{karpukhin2020}
KARPUKHIN V, OGUZ B, MIN S, et al. Dense Passage Retrieval for Open-Domain Question Answering[C]//Proceedings of the 2020 Conference on Empirical Methods in Natural Language Processing. Stroudsburg, PA: Association for Computational Linguistics, 2020: 6769-6781. DOI:10.18653/v1/2020.emnlp-main.550.

\bibitem{xiao2024}
XIAO S, LIU Z, ZHANG P, et al. C-Pack: Packed Resources for General Chinese Embeddings[C]//Proceedings of the 47th International ACM SIGIR Conference on Research and Development in Information Retrieval. New York: ACM, 2024: 641-649. DOI:10.1145/3626772.3657878.

\bibitem{gao2025}
郜洪奎, 马瑞祥, 包骐豪, 等. 基于混合检索增强的双塔模型研究[J]. 计算机科学, 2025, 52(6): 324-329. DOI:10.11896/jsjkx.240800017.

\bibitem{zhangjian2025}
张健, 唐晋韬, 王挺, 等. 基于混合检索重排序策略的大模型增强方法[J]. 中文信息学报, 2025, 39(4): 42-54.

\bibitem{sun2026}
孙飒, 张经伟, 江卫. 面向企业知识库的层次化分块与混合检索增强生成方法设计[J]. 通信技术, 2026, 59(7): 745-754. DOI:10.3969/j.issn.1002-0802.2026.07.002.

\bibitem{song2026}
宋保仁, 冷静. 检索增强生成技术在图书馆私有知识库构建中的应用与边界——以四川高校特色数据库为例[J]. 情报理论与实践, 2026, 49(4): 184-193. DOI:10.16353/j.cnki.1000-7490.2026.04.021.

\bibitem{qin2026}
秦晓瑞, 何柳, 安然, 等. 面向航空标准的大语言模型迭代检索增强生成方法[J]. 南京航空航天大学学报（自然科学版）, 2026, 58(1): 235-248. DOI:10.16356/j.2097-6771.2026.01.022.

\bibitem{gan2026}
甘克勤, 高亮, 肖宝坤, 等. 基于检索增强生成与对话历史管理的标准智能问答系统研究[J]. 中国标准化, 2026(2): 43-47. DOI:10.3969/j.issn.1002-5944.2026.02.002.

\bibitem{luo2026}
罗云, 徐嘉骏, 彭明阳, 等. 忠实性增强的国防政策检索增强生成方法[J]. 火箭军工程大学学报, 2026, 40(1): 132-142. DOI:10.20189/j.cnki.CN/61-1527/E.202601012.

\bibitem{dong2025}
董永涛, 崔翛龙. 基于大模型和检索增强生成技术的舰艇装备故障智能问答系统[J]. 中国舰船研究, 2025, 20(6): 144-153. DOI:10.19693/j.issn.1673-3185.04285.

\bibitem{peng2026}
彭琳. 基于检索增强生成的企业级知识库系统构建与优化[J]. 电脑知识与技术, 2026, 22(12): 35-37.

\bibitem{tao2026}
陶金涛, 张楠楠, 廖世斌. 基于检索增强生成与大语言模型的新疆地质矿产智能问答系统构建研究[J]. 新疆地质, 2026, 44(1): 154-159.

\bibitem{wang2026}
王祥昊, 王莉, 胡恒闯, 等. 基于知识检索增强的铁路应急处置策略智能生成方法[J]. 中国铁道科学, 2026, 47(1): 233-243. DOI:10.3969/j.issn.1001-4632.2026.01.21.

\bibitem{dou2026}
窦嘉铭, 王小君, 乔骥, 等. 基于知识图谱检索增强生成与大语言模型的配网检修智能辅助决策框架[J]. 工程科学与技术, 2026, 58(4): 67-79.

\bibitem{es2024}
ES S, JAMES J, ESPINOSA-ANKE L, et al. RAGAs: Automated Evaluation of Retrieval Augmented Generation[C]//Proceedings of the 18th Conference of the European Chapter of the Association for Computational Linguistics: System Demonstrations. Stroudsburg, PA: Association for Computational Linguistics, 2024: 150-158. DOI:10.18653/v1/2024.eacl-demo.16.

\bibitem{lu2025}
卢万有. 基于JWT的RBAC在前后端分离项目中的设计与实现[J]. 电脑编程技巧与维护, 2025(1): 46-48.

\bibitem{liao2026}
廖明焱. 面向企业应用的RAG知识安全检索与AI脱敏技术研究[J]. 电脑编程技巧与维护, 2026(8): 126-128.

\bibitem{wei2025}
魏巍. 基于RBAC模型云资源管理系统访问控制权限的设计与实施[J]. 科学技术创新, 2025(2): 122-125.

\bibitem{zhangning2026}
张宁. 基于LangChain的Python知识问答系统的设计与实现[J]. 现代信息科技, 2026, 10(5): 66-71. DOI:10.19850/j.cnki.2096-4706.2026.05.013.

\bibitem{wu2026}
吴海周, 黄俊, 张佳晨. 基于RAG机制的知识库问答性能优化与参数调控研究[J]. 工业控制计算机, 2026, 39(7): 111-113.

\bibitem{langchain2026}
LangChain. LangChain Documentation[EB/OL]. [2026-09-21]. \url{https://python.langchain.com/docs/}.

\end{thebibliography}
